\section{Two quantities that bound the maps}\label{sec:theory}
\paragraph{Population versus context-only: a fixed budget.} Let $P$ and $Q$ be the population and context-only priors over a
system $\theta$, with the same likelihood, and let $g_n$ be the expected one-step gap at context length $n$ on held-out
systems ($\theta\sim P$). Write $m_P^N$ and $m_Q^N$ for the two Bayes marginal laws of $x_{1:N}$ given $x_0$; because $x_0$ is
drawn independently of $\theta$, the one-step predictives of the two references are the conditionals of these marginals. By the
chain rule of relative entropy \citep{cover2006},
\[
\textstyle\sum_{n=0}^{N-1} g_n \;=\; D\!\left(m_P^N \,\middle\|\, m_Q^N\right) \;\le\; D(P\,\|\,Q),
\]
where the inequality holds because $m_P^N$ and $m_Q^N$ are marginals of the joint laws of $(\theta,x_{1:N})$, which share the
likelihood and differ only in the prior, and marginalising cannot increase relative entropy. The left side is the cumulative
log loss (redundancy) of forecasting with $Q$'s mixture when the data come from $P$'s; for a single prior, the analogous
cumulative risk of the Bayes strategy is a mutual information \citep{haussler1997}. The bound is classical; we use it to
predict, before any map, how much room a family leaves. It has three consequences. The per-forecast gaps of a family spend a
fixed budget $D(P\|Q)$; contexts that are informative early (large initial states) spend it early, so their window is
short. Observation noise known to both references leaves the budget unchanged and spreads it over more forecasts, so the
per-forecast gaps shrink while the total population advantage does not. And because the priors in F1 are products over the
coefficient and the noise scale, the budget splits exactly into a dynamics part and a noise-level part.

The budgets (relative entropy between priors, computed numerically) are \KLoneA nats with a known noise scale (and in F3,
whose priors over $a$ are the same); \KLFone nats in F1, of which \KLFoneSigma come from the noise-scale priors, a part set by
the width of the context-only log-uniform range (\KLFoneSigmaNarrow nats if it spanned two orders of magnitude instead of
four); and \KLFtwo nats in
F2 (\KLFtwoNM if the context-only reference were given the population's noise prior). So most of the room in F1 and F2 is
knowledge of the noise level, which any reference trained on the population's data would acquire.

\paragraph{Library versus population: the nearest system.} On held-out systems, as the context grows, the population
predictive approaches the true one-step law and the library posterior concentrates on the library system nearest to the
truth in per-step divergence, so the gap tends to $\mathbb E[\min_k D_{\mathrm{stat}}(\theta\,\|\,\theta_k)]$, the expected
per-step divergence under the stationary law of the regressor (a derivation that assumes posterior concentration; finite
contexts lie before this limit). Among $K$ systems drawn from a $p$-parameter population the nearest lies at a distance of
order $K^{-1/p}$, so the limit falls roughly like $K^{-2/p}$: fast in scalar families and slowly for the seven-parameter
bivariate systems. By Monte Carlo, the limit is \NaKnownFour, \NaKnownEight and \NaKnownSixteen nats for libraries of
$K=4$, $8$ and \KSixteen systems with a known noise scale; \NaUnknownSixteen and \NaUnknownThirtyTwo nats for $K=\KSixteen$
and \KThirtyTwo when the noise scale is unknown; and \NaTwoSixtyFour nats in the bivariate family at $K=\FtwoKmax$, still
\NaTwoThousand at $K=\KThousand$. At $\delta=\DeltaNats$, the limit therefore falls below the margin between four and eight
systems with a known noise scale, between \KSixteen and \KThirtyTwo with an unknown one, and not within the library sizes we
evaluated in the bivariate family.
